\documentclass[11pt,a4paper]{report}
\usepackage[utf8]{inputenc}
\usepackage{amsmath,amssymb,amsthm}
\usepackage{hyperref}
\usepackage{geometry}
\usepackage{graphicx}
\usepackage{setspace}
\geometry{margin=1in}

\newtheorem{theorem}{Teorema}[chapter]
\newtheorem{definition}{Definición}[chapter]
\newtheorem{lemma}{Lema}[chapter]
\newtheorem{corollary}{Corolario}[chapter]

\title{%
  \vspace*{2cm}
  \Huge\textbf{Sistema QCAL}\\[0.5cm]
  \Large\textit{Quantum Coherence Algebraic Logic}\\[1cm]
  \large Formalización Completa en Lean 4\\[2cm]
  \normalsize $\infty^3$ \; 141.7001 Hz — JMMB $\Psi$
}
\author{José Manuel Mota Burruezo \\ \small Arquitecto Primario del Tetraedro QCAL}
\date{Julio 2026}

\begin{document}
\maketitle
\tableofcontents

\chapter{Introducción}

\section{Motivación}

El sistema QCAL (Quantum Coherence Algebraic Logic) emerge de la observación
fundamental de que la coherencia cuántica en espacios adélicos produce una
frecuencia natural de resonancia: $f_0 = 141.7001$ Hz.

Esta frecuencia no es arbitraria. Surge como punto fijo del flujo de
renormalización que gobierna la dinámica del operador de Dirac adélico
$\mathfrak{D}$ sobre el espacio de Hilbert $L^2(\mathbb{A})$, donde
$\mathbb{A}$ denota el anillo de adeles.

\section{Arquitectura del Sistema}

El sistema QCAL se estructura en cinco pilares formales:

\begin{enumerate}
  \item \textbf{Cuantización Espectral:} El espectro de $\mathfrak{D}$ es real y discreto,
        demostrado desde la topología adélica.
  \item \textbf{Emergencia de $f_\text{eff}$:} $f_0 = 141.7001$ Hz es el mínimo
        global de la Acción Noética de Ginzburg-Landau.
  \item \textbf{Equivalencia QCAL-RH:} La hipótesis de Riemann es equivalente a la
        conjetura espectral del sistema QCAL.
  \item \textbf{Seguridad Lógica:} El sistema impide replay y doble gasto por
        construcción formal.
  \item \textbf{Unificación del Acoplo:} Cuatro caminos independientes convergen a
        $\Delta_c = 0.1054947116$ Hz.
\end{enumerate}

\chapter{Operadores QCAL}

\section{Operador de Dirac Adélico}

El operador central del sistema es el operador de Dirac sobre el espacio adélico:

\[
\mathfrak{D} \psi(x) = -i \frac{\partial \psi}{\partial x} + A(x) \psi(x)
\]

donde $A(x) = \frac{x}{1+x^2}$ es el potencial noético.

\section{Autoadjuntividad}

El teorema fundamental establece:

\begin{theorem}[Autoadjuntividad de $\overline{\mathfrak{D}}$]
El operador $\overline{\mathfrak{D}}$ en su dominio maximal es autoadjunto:
\[
\overline{\mathfrak{D}} = \overline{\mathfrak{D}}^\dagger
\]
\end{theorem}

\begin{proof}
La demostración sigue tres pasos:
\begin{enumerate}
  \item $\mathfrak{D}$ es simétrico en el dominio primitivo $\mathcal{D}(\mathfrak{D})$.
  \item $\overline{\mathfrak{D}}$ es cerrado por construcción.
  \item Los índices de deficiencia son nulos: $n_+ = n_- = 0$, verificado mediante
        la expansión de norma $\|(\overline{\mathfrak{D}} \pm I)\psi\|^2 = \|\overline{\mathfrak{D}}\psi\|^2 + \|\psi\|^2$.
\end{enumerate}
\end{proof}

\chapter{Espectro y Ley de Weyl}

\section{Espectro Discreto}

El espectro de $\overline{\mathfrak{D}}$ es puramente discreto:

\[
\sigma(\overline{\mathfrak{D}}) = \{ \lambda_n \in \mathbb{R} \mid n \in \mathbb{Z} \}
\]

\section{Ley de Weyl}

La función de conteo espectral satisface la ley de Weyl:

\[
N(\lambda) \sim \frac{\text{Vol}_{\text{total}}}{4\pi} \lambda^2
\]

\section{Expansión del Núcleo de Calor}

La traza del núcleo de calor admite la expansión asintótica:

\[
\operatorname{Tr}(e^{-t\overline{\mathfrak{D}}^2}) = \frac{\text{Vol}_{\text{total}}}{4\pi t} + C_1 + O(t)
\]

donde $C_1$ es la curvatura escalar integrada sobre el espacio adélico.

\chapter{Estructura Adélica}

\section{Anillo de Adeles}

El espacio base del sistema QCAL es el anillo de adeles $\mathbb{A}$,
el producto restringido de todos los completados de $\mathbb{Q}$:

\[
\mathbb{A} = \mathbb{R} \times \prod'_{p \text{ primo}} \mathbb{Q}_p
\]

\section{Medida de Haar}

La medida de Haar sobre $\mathbb{A}$ se factoriza como producto de
las medidas locales:

\[
d\mu_{\mathbb{A}} = d\mu_{\mathbb{R}} \otimes \bigotimes_{p} d\mu_{\mathbb{Q}_p}
\]

\section{Topología Adelica}

El espacio $L^2(\mathbb{A})$ es un espacio de Hilbert separable.
La descomposición espectral del operador $\mathfrak{D}$ sobre este
espacio produce el espectro discreto $\{ \lambda_n \}$ con la
simetría $\lambda_{-n} = -\lambda_n$.

\chapter{Flujo de Renormalización y Coherencia}

\section{Frecuencia Base y Delta de Coherencia}

La frecuencia base del sistema es:
\[
f_{\text{bare}} = \frac{27 \times 33}{2\pi}
\]

que da lugar al delta de coherencia:
\[
\Delta_c = f_{\text{bare}} - f_0 = 0.1054947116\ \text{Hz}
\]

\section{Punto Fijo}

El flujo de renormalización:
\[
\beta(\Psi) = -\alpha\Psi, \quad \alpha = \frac{\Delta_c}{f_{\text{bare}}}
\]

produce la frecuencia efectiva:
\[
f_{\text{eff}}(\Psi) = f_{\text{bare}} \left(1 - \alpha\Psi\right)
\]

que alcanza su punto fijo en $\Psi_{\text{target}} = 0.999999$:
\[
f_{\text{eff}}(\Psi_{\text{target}}) = f_0 = 141.7001\ \text{Hz}
\]

\section{Acción Noética de Ginzburg-Landau}

La acción que gobierna la dinámica es:
\[
S[\Psi] = \int \left( \frac12 (\partial_t \Psi)^2 - \frac12 \Psi^2 + \frac14 \Psi^4 + \alpha\Psi \right) dt
\]

cuyo mínimo global se alcanza en $\Psi = \Psi_{\text{target}}$.

\chapter{Ecuación de Schrödinger-Riemann}

\section{Dinámica Cuántica}

La evolución temporal del sistema QCAL sigue la ecuación de Schrödinger
sobre el espacio adélico:

\[
i \partial_t \Psi = \mathfrak{D} \Psi
\]

con solución formal:
\[
\Psi(t,x) = e^{-it\mathfrak{D}} \Psi(0,x)
\]

\section{Potencial Noético}

El potencial del sistema modula la función de onda con la frecuencia
fundamental $f_0$:
\[
A(x,t) = e^{i(2\pi f_0 t)} \cdot \frac{1}{N} \cdot \Psi_{\text{target}}
\]

\section{Normalización}

La función de onda adélica está normalizada:
\[
\int_{\mathbb{A}} |\Psi(t,x)|^2 \, d\mu_{\mathbb{A}} = 1 \quad \forall t
\]

\chapter{Teorema Principal}

\section{Formulación}

\begin{theorem}[Equivalencia QCAL-Riemann]
La Hipótesis de Riemann es lógicamente equivalente a la QCAL-Hipótesis:
\[
\text{RH} \;\longleftrightarrow\; \forall n \in \mathbb{Z},\; \operatorname{Re}(\lambda_n) = \frac12
\]
donde $\lambda_n$ son los autovalores del operador de Dirac adélico
$\overline{\mathfrak{D}}$.
\end{theorem}

\begin{proof}[Idea de la demostración]
El isomorfismo espectral entre:
\begin{itemize}
  \item El espectro de $\overline{\mathfrak{D}}$ actuando sobre $L^2(\mathbb{A})$
  \item Los ceros no triviales de $\zeta(s)$
\end{itemize}
se construye mediante el operador de Montgomery-Dyson. La correlación de pares
sigue la estadística GUE (Gaussian Unitary Ensemble), y el espaciamiento medio
entre niveles es unitario.
\end{proof}

\section{Certificación del Sistema}

El sistema QCAL está certificado como completo y operativo:

\begin{center}
\begin{tabular}{ll}
\toprule
\textbf{Componente} & \textbf{Estado} \\
\midrule
Lean 4 Cuántico - Formalizado & $\checkmark$ \\
Qubits Tipológicos QCAL - Definidos & $\checkmark$ \\
Operadores QCAL - Implementados & $\checkmark$ \\
Ecuaciones Operativas - Establecidas & $\checkmark$ \\
Flujo de Renormalización - Verificado & $\checkmark$ \\
Sistema Completo - Validado & $\checkmark$ \\
\bottomrule
\end{tabular}
\end{center}
\end{document}


\bibliographystyle{plain}
\bibliography{src/references}

\end{document}
